%% ch06.tex — Chapter 6: Apache Spark: In-Memory Distributed Computing
%% Part of "Data Engineering" textbook

\chapter{Apache Spark: In-Memory Distributed Computing}
\label{ch:spark}

\epigraph{\textit{``The most important question we faced was: what should be
the fundamental programming model? We chose to expose distributed memory
as a first-class abstraction, and that single decision determined
everything else about Spark's design.''}}{Matei Zaharia, co-creator of Apache Spark, 2016}

\vspace{4pt}

\begin{chapterlearning}
After completing this chapter, you will be able to:
\begin{enumerate}
  \item Identify the specific I/O bottleneck of MapReduce for iterative
        workloads and derive the disk-I/O cost formula that motivated
        the design of Spark.
  \item Define a Resilient Distributed Dataset (RDD), explain its five
        defining properties, and describe how lineage-based fault tolerance
        differs from replication-based fault tolerance.
  \item Trace the life cycle of a Spark job from \texttt{sc.textFile()}
        through the DAG Scheduler, Task Scheduler, and Executor, identifying
        the role of each component.
  \item Distinguish narrow transformations from wide transformations,
        explain why wide transformations introduce shuffle boundaries,
        and identify which RDD operations fall into each category.
  \item Explain the performance difference between \texttt{groupByKey}
        and \texttt{reduceByKey}, and apply the correct operator to avoid
        unnecessary data movement.
  \item Describe the four phases of the Catalyst query optimizer (analysis,
        logical optimisation, physical planning, and code generation) and
        explain how each phase reduces query execution cost.
  \item Select the appropriate RDD persistence level from
        \texttt{MEMORY\_ONLY}, \texttt{MEMORY\_AND\_DISK},
        \texttt{MEMORY\_ONLY\_SER}, and \texttt{DISK\_ONLY} for a given
        workload, cost, and fault-tolerance requirement.
  \item Describe the Spark MLlib Pipeline API---Transformer, Estimator,
        Pipeline, and PipelineModel---and construct a three-stage
        feature engineering and classification pipeline.
  \item Explain Spark's unified memory management model, distinguish
        execution memory from storage memory, and calculate the available
        memory for each region given a heap size and configuration.
  \item Apply the salting technique to resolve data skew in a wide
        transformation, and explain why skew is undetectable by the
        Catalyst optimizer.
\end{enumerate}
\end{chapterlearning}

\bigskip

By 2010, Hadoop's MapReduce framework had conquered the storage and
batch-processing challenges that motivated its creation. Thousands of
organisations processed petabytes per day with clusters of commodity
servers, achieving the economics that Google had demonstrated internally
with its MapReduce implementation \parencite{MapReduce2004}. Yet
practitioners who attempted to use MapReduce for machine learning training,
graph analysis, and iterative data mining encountered a fundamental
performance barrier: MapReduce was designed for data that is read once,
processed once, and written once. Any algorithm that reads the same data
repeatedly---which describes virtually every machine learning training
algorithm---paid an enormous, repeated disk I/O penalty that made
computation on even modest datasets prohibitively slow.

This chapter introduces Apache Spark, the distributed computing framework
that resolved the iterative-workload bottleneck by elevating distributed
memory to a first-class programming abstraction. Spark's central concept,
the \term{Resilient Distributed Dataset} (RDD), allows the results of an
expensive computation to be stored in the combined RAM of a cluster and
reused across multiple algorithm iterations without a single disk read.
For workloads such as logistic regression training, this in-memory
approach achieved speedups of 20$\times$ to 100$\times$ over MapReduce
in the original benchmarks \parencite{RDD2012}.

%%====================================================================
\section{The MapReduce Iteration Problem}
\label{sec:mr-iteration}
%%====================================================================

To understand why Spark was necessary, we must quantify precisely what
MapReduce does poorly: iterative computation over the same dataset.

Consider training a logistic regression model on a 100\GB\ training set
using gradient descent. The training algorithm proceeds in rounds, called
\emph{iterations} or \emph{epochs}. In each iteration, the algorithm:
(1) reads the entire training set; (2) computes the gradient of the loss
function with respect to the current model parameters; (3) updates the
parameters; and (4) checks a convergence criterion. A typical training
run requires 50 to 200 iterations before the model converges.

\textbf{The MapReduce implementation} maps one iteration to one
MapReduce job. Each job:
\begin{enumerate}
  \item Reads the 100\GB\ training set from HDFS (3$\times$ replicated
        = 300\GB\ of disk reads).
  \item Runs Map tasks to compute per-record gradient contributions.
  \item Runs a Reduce task to aggregate the gradients.
  \item Writes the updated model parameters to HDFS.
  \item The next iteration reads HDFS again from the beginning.
\end{enumerate}

\begin{defbox}{MapReduce Iteration I/O Cost}
For an iterative algorithm requiring $I$ iterations over a dataset of
size $D$ bytes, stored on HDFS with replication factor $r$:
\begin{align}
\text{Total disk I/O (MapReduce)} &= I \times D \times r
  \label{eq:mr-io-cost}
\end{align}
For logistic regression with $I = 100$ iterations, $D = 100\GB$,
$r = 3$:
\[
  \text{Total disk I/O} = 100 \times 100\GB \times 3 = 30{,}000\GB = 30\TB
\]
Processing 30\TB\ of disk I/O to train a model on 100\GB\ of unique data
is a $300\times$ I/O amplification. On a cluster that sustains 1\GB/s
aggregate disk read throughput per node, with 100 nodes, this represents
5 minutes of wall-clock time \emph{per iteration}---8.3 hours total for
100 iterations.
\end{defbox}

The root cause is straightforward: MapReduce has no mechanism for retaining
intermediate data in memory between jobs. Every job boundary is a forced
write to and read from HDFS. The framework was designed for single-pass
batch processing, not iterative refinement. The cluster's RAM---which by
2010 already totalled hundreds of gigabytes across even a modest 100-node
cluster---is completely wasted between iterations: it is cleared, and the
entire dataset is re-read from disk for the next round.

Interactive analytical workloads suffer a related problem. A data
analyst running ten queries over the same dataset against a Hive-on-MapReduce
cluster must wait for the full input to be read from HDFS for each query,
even if the dataset has already been read by a prior query. There is no
shared memory pool between independent jobs that could serve as a cache.

Matei Zaharia and his collaborators at UC Berkeley identified both problems
in 2010 and proposed a solution: expose the cluster's aggregate RAM as a
distributed, fault-tolerant storage layer that persists across job
boundaries. This idea became Apache Spark \parencite{Spark2010}.

%%====================================================================
\section{Resilient Distributed Datasets}
\label{sec:rdd}
%%====================================================================

The fundamental abstraction of Apache Spark is the \term{Resilient
Distributed Dataset} (RDD), introduced formally by Zaharia et al.\ at
NSDI 2012 \parencite{RDD2012}. An RDD is the unit of data in Spark in
the same way that an HDFS block is the unit of data in Hadoop: every
computation operates on RDDs, produces RDDs, and stores RDDs.

\begin{defbox}{Resilient Distributed Dataset (RDD)}
An \textbf{RDD} is an immutable, partitioned, fault-tolerant collection
of records that can be stored in memory across a cluster of nodes and
computed in parallel. Formally, an RDD $R$ is characterised by five
properties:
\begin{enumerate}
  \item \textbf{A set of partitions:} $R = \{P_1, P_2, \ldots, P_k\}$,
        where each partition $P_i$ is a subset of the data assigned to
        one executor.
  \item \textbf{A set of dependencies:} a reference to the parent RDD(s)
        from which $R$ was derived, forming a lineage graph (DAG).
  \item \textbf{A function to compute each partition} from its parent
        partition(s).
  \item \textbf{Partitioning metadata:} the scheme (e.g., hash or range)
        used to assign records to partitions.
  \item \textbf{Preferred locations:} hints on which nodes hold the data
        locally (e.g., HDFS block locations).
\end{enumerate}
\end{defbox}

%--------------------------------------------------------------------
\subsection{Lineage-Based Fault Tolerance}
\label{subsec:lineage}
%--------------------------------------------------------------------

The ``Resilient'' in RDD refers to Spark's mechanism for recovering from
executor failures: \term{lineage-based fault tolerance}. This is the most
conceptually important property of RDDs, and it distinguishes Spark's
approach from both MapReduce (which writes intermediate results to disk)
and traditional distributed shared memory (which replicates data in RAM).

When an executor fails and the partitions it held are lost, Spark does not
need to restart the entire job. Instead, it consults the RDD's
\term{lineage graph}---the directed acyclic graph of transformations that
produced the lost partition---and \emph{recomputes only the lost partition}
on a surviving executor.

\begin{figure}[h!t]
\centering
\begin{tikzpicture}[
    rdd/.style={rectangle, draw=DEBlue!60, fill=DELightBlue!50,
                rounded corners=5pt, minimum width=2.8cm, minimum height=0.7cm,
                font=\small\bfseries\sffamily, align=center},
    part/.style={rectangle, draw=DEGreen!60, fill=DELightGreen,
                 rounded corners=3pt, minimum width=1.1cm, minimum height=0.55cm,
                 font=\tiny\sffamily, align=center},
    lostpart/.style={rectangle, draw=DERed!70, fill=DERed!20,
                     rounded corners=3pt, minimum width=1.1cm, minimum height=0.55cm,
                     font=\tiny\sffamily, align=center, dashed},
    arr/.style={-{Stealth[length=4pt,width=4pt]}, thick},
    lbl/.style={font=\scriptsize\sffamily, align=center}
  ]

  % RDD A (source)
  \node[rdd] (A)  at (0,   5.2) {RDD A\\(source)};
  \node[part] (A1) at (-1.5, 4.2) {$P_1$};
  \node[part] (A2) at (0,   4.2) {$P_2$};
  \node[part] (A3) at (1.5, 4.2) {$P_3$};
  \draw[arr, DEBlue!50] (A.south) -- (A1.north);
  \draw[arr, DEBlue!50] (A.south) -- (A2.north);
  \draw[arr, DEBlue!50] (A.south) -- (A3.north);

  % Transformation label
  \node[lbl, text=DEAccent] at (3.5, 4.5) {.filter(f)\\(narrow dep.)};
  \draw[DEAccent!50, dashed, -{Stealth[length=3pt]}] (3.5, 4.1) -- (1.0, 4.1);

  % RDD B (derived)
  \node[rdd] (B)  at (0, 3.0) {RDD B\\(filtered)};
  \node[part]     (B1) at (-1.5, 2.0) {$P_1$};
  \node[lostpart] (B2) at (0,   2.0) {$P_2$ \textbf{LOST}};
  \node[part]     (B3) at (1.5, 2.0) {$P_3$};
  \draw[arr, DEBlue!50] (B.south) -- (B1.north);
  \draw[arr, DERed!60,  dashed]   (B.south) -- (B2.north);
  \draw[arr, DEBlue!50] (B.south) -- (B3.north);

  % Lineage arrows from A partitions to B partitions
  \draw[arr, DEGreen!70!black, thick] (A1.south) -- (B1.north);
  \draw[arr, DERed!60, dashed, thick] (A2.south) -- (B2.north);
  \draw[arr, DEGreen!70!black, thick] (A3.south) -- (B3.north);

  % Recovery annotation
  \node[lbl, text=DERed, align=left] at (4.5, 2.3)
       {\textbf{Executor fails:}\\$P_2$ of RDD~B lost};
  \node[lbl, text=DEGreen!70!black, align=left] at (4.5, 1.4)
       {\textbf{Recompute:}\\reapply \texttt{filter(f)}\\to $A.P_2$ only};
  \draw[DERed, -{Stealth[length=3pt]}] (3.4, 2.1) -- (0.6, 2.1);
  \draw[DEGreen!70!black, -{Stealth[length=3pt]}] (3.4, 1.5) -- (0.6, 1.5);

  % Lineage label
  \node[font=\small\bfseries\sffamily, text=DEGray] at (-4.0, 3.5)
       {Lineage};
  \node[font=\tiny\sffamily, text=DEGray, align=center] at (-4.0, 3.1)
       {``$B$ was computed\\by applying filter~$f$\\to each partition of~$A$''};
  \draw[DEGray!50, dashed] (-3.2, 3.5) -- (-1.7, 3.5);

\end{tikzpicture}
\caption{RDD lineage-based fault tolerance. When $P_2$ of RDD~B is lost
         due to executor failure, Spark consults the lineage graph to determine
         that $P_2$ of~B was produced by applying \texttt{filter(f)} to $P_2$
         of RDD~A, which is still available. Only the lost partition is
         recomputed---no full job restart is required.}
\label{fig:rdd-lineage}
\end{figure}

This design trades replication overhead for recomputation cost. RDDs
backed by data on HDFS (the most common case) can always recompute any
lost partition from its HDFS source data, which is itself replicated.
For RDDs produced by long chains of transformations, Spark allows
users to insert \term{checkpoints}---explicit writes of an RDD to HDFS
that truncate the lineage graph---preventing unbounded recomputation
cost in the event of failure after many transformation stages.

\begin{keyinsight}{Lineage vs.\ Replication: A Design Trade-off}
Lineage-based fault tolerance costs nothing in memory (no second copy
of the data is maintained) but can be expensive in time if the lineage
is long. Replication-based fault tolerance (as in HDFS) costs memory
for every replica but recovers in constant time regardless of lineage
length. Spark's design assumes that executor failures are rare and
datasets are large relative to available RAM---so paying memory for
replication would defeat the purpose of in-memory computing. For
streaming workloads where maintaining lineage across an unbounded stream
is impractical, Spark Structured Streaming adds micro-batch replication.
\end{keyinsight}

%--------------------------------------------------------------------
\subsection{Lazy Evaluation and the RDD DAG}
\label{subsec:lazy}
%--------------------------------------------------------------------

RDDs are computed \term{lazily}: calling a transformation method such as
\texttt{filter()} or \texttt{map()} does not immediately execute any
computation. It only records the transformation in the RDD's lineage
graph. Actual computation is triggered only when the programmer calls
an \term{action}---a method that returns a non-RDD result to the driver
(such as \texttt{count()}, \texttt{collect()}, or \texttt{saveAsTextFile()}).

This laziness is not a limitation; it is a deliberate optimisation
opportunity. When an action is called, Spark's DAG Scheduler can inspect
the entire lineage graph of all RDDs involved in producing the result
and optimise the computation plan. It can fuse adjacent narrow
transformations into a single pass over the data, avoiding intermediate
materialisation, and it can identify which partitions are already cached
in memory and skip their recomputation.

%--------------------------------------------------------------------
\subsection{RDD Persistence}
\label{subsec:persistence}
%--------------------------------------------------------------------

The performance advantage of Spark over MapReduce materialises only when
RDDs are \emph{persisted}. When an RDD is used in multiple actions---the
common case in iterative algorithms---the programmer must explicitly call
\texttt{rdd.persist(StorageLevel.MEMORY\_ONLY)} (or the equivalent
\texttt{rdd.cache()}, which is shorthand for
\texttt{MEMORY\_ONLY}) to instruct Spark to retain the RDD's partitions
in executor memory after the first computation. Without persistence,
each action triggers a full recomputation from the lineage root.

\begin{table}[h!t]
\centering
\caption{RDD storage levels and their trade-offs}
\label{tab:storage-levels}
\renewcommand{\arraystretch}{1.5}
\begin{tabular}{L{3.2cm} L{1.4cm} L{1.4cm} L{4.5cm}}
\toprule
\textbf{Storage level} & \textbf{Memory} & \textbf{Disk} & \textbf{When to use} \\
\midrule
\texttt{MEMORY\_ONLY}       & Yes (raw) & No  & Fits in RAM; fast recompute from HDFS \\
\texttt{MEMORY\_AND\_DISK}  & Yes (raw) & Overflow & Larger-than-RAM; avoid full recompute \\
\texttt{MEMORY\_ONLY\_SER}  & Yes (Kryo) & No & Tight on RAM; CPU overhead acceptable \\
\texttt{MEMORY\_AND\_DISK\_SER} & Yes (Kryo) & Overflow & Tight on RAM + large data \\
\texttt{DISK\_ONLY}         & No        & Yes & Rarely accessed; fast recompute is slow \\
\texttt{OFF\_HEAP}          & Off-heap  & No  & Avoid GC pauses; Tungsten memory \\
\bottomrule
\end{tabular}
\end{table}

Spark evicts persisted RDD partitions using a \emph{least-recently-used}
(LRU) policy when storage memory is exhausted. Under
\texttt{MEMORY\_AND\_DISK}, evicted partitions spill to the local disk of
the executor node rather than being dropped; they are read back from disk
if needed by a subsequent action. Under \texttt{MEMORY\_ONLY}, evicted
partitions are simply recomputed from their lineage on demand.

%%====================================================================
\section{Spark Architecture}
\label{sec:spark-arch}
%%====================================================================

A Spark application consists of a single \term{driver program} and one
or more \term{executors} running on worker nodes. The driver and
executors communicate through a \term{cluster manager}, which allocates
the physical resources (CPU cores and RAM) that the application requires.

%--------------------------------------------------------------------
\subsection{The Driver Program and SparkSession}
\label{subsec:driver}
%--------------------------------------------------------------------

The \term{driver program} is the process that contains the user's
application code. It runs the \texttt{main()} function of the Spark
application, constructs the \term{SparkSession} (the unified entry point
since Spark 2.0, replacing the earlier \texttt{SparkContext} and
\texttt{HiveContext}), and orchestrates the overall execution.

The driver performs three critical functions:
\begin{itemize}
  \item \textbf{DAG construction:} as the application code calls
        transformation methods, the driver records each transformation
        in an RDD lineage DAG. No computation occurs yet.
  \item \textbf{Job submission:} when an action is called, the driver
        submits a job to the DAG Scheduler. The DAG Scheduler partitions
        the lineage DAG into \term{stages}---groups of transformations
        that can be pipelined without a data shuffle---and submits each
        stage as a set of \term{tasks} to the Task Scheduler.
  \item \textbf{Result collection:} for actions that return data to the
        driver (such as \texttt{collect()} or \texttt{first()}), the
        driver receives task results from executors and assembles the
        final result.
\end{itemize}

\begin{caution}{The Driver Memory Trap}
\texttt{rdd.collect()} transfers \emph{all} partitions of an RDD to the
driver's JVM heap. On a production dataset of even a few hundred
gigabytes, this causes an OutOfMemoryError that crashes the driver
process and terminates the entire application. The correct alternatives
are \texttt{rdd.take(n)} (return only the first $n$ records),
\texttt{rdd.saveAsTextFile(path)} (write results to HDFS without
collecting), or \texttt{rdd.sample(false, 0.01)} (return a 1\% sample).
Use \texttt{collect()} only for small result sets that the driver
can safely hold in memory---such as the output of a final aggregation.
\end{caution}

%--------------------------------------------------------------------
\subsection{Executors}
\label{subsec:executors}
%--------------------------------------------------------------------

\term{Executors} are long-lived JVM processes launched on worker nodes
at application start-up time. Each executor:
\begin{itemize}
  \item Runs \emph{tasks} assigned by the Task Scheduler, with multiple
        tasks running concurrently in separate threads (one thread per
        available vCPU core allocated to the executor).
  \item Maintains a local \emph{block store}---the Spark Block Manager---
        that holds cached RDD partitions and shuffle output blocks.
  \item Reports task status, progress, and shuffle data locations back
        to the driver via a heartbeat mechanism.
  \item Runs for the entire lifetime of the application, enabling in-memory
        data sharing across multiple Spark jobs within the same application.
\end{itemize}

The executor's JVM heap is divided between \emph{execution memory}
(for active task computation: sort buffers, hash tables for aggregations,
shuffle read/write buffers) and \emph{storage memory} (for cached
RDD partitions). This division is discussed in depth in
Section~\ref{sec:memory}.

%--------------------------------------------------------------------
\subsection{Cluster Managers}
\label{subsec:cluster-managers}
%--------------------------------------------------------------------

Spark is \emph{cluster-manager-agnostic}: it delegates the allocation
of physical resources to a pluggable cluster manager. Three cluster
managers are commonly used in production:

\begin{itemize}
  \item \textbf{Spark Standalone:} Spark's built-in cluster manager,
        suitable for dedicated Spark clusters. Simple to configure but
        does not support multi-tenant resource sharing or integration
        with other frameworks. Appropriate for development and
        single-purpose deployment.
  \item \textbf{YARN:} the most common production deployment (see
        Chapter~\ref{ch:yarn-formats}). Spark submits an
        ApplicationMaster to YARN; the AM negotiates executor containers
        with the ResourceManager. All of YARN's scheduling policies
        (Capacity, Fair, DRF) apply to Spark executors. Enables
        multi-framework resource sharing with MapReduce, Tez, and Samza
        on the same cluster.
  \item \textbf{Kubernetes:} the preferred deployment model on
        cloud infrastructure. The driver runs as a Kubernetes Pod; each
        executor runs in its own Pod. Kubernetes manages container
        scheduling, scaling, and isolation using its standard
        infrastructure, enabling integration with cloud autoscaling.
\end{itemize}

%--------------------------------------------------------------------
\subsection{DAG Scheduler and Task Scheduler}
\label{subsec:dag-task-sched}
%--------------------------------------------------------------------

When an action triggers job execution, the driver's scheduler pipeline
proceeds in two phases.

The \term{DAG Scheduler} receives the RDD lineage graph, identifies the
\emph{stage boundaries} (discussed in Section~\ref{sec:transformations}),
and decomposes the computation into a DAG of stages. Each stage is a
maximal set of consecutive narrow transformations that can be pipelined.
Stage boundaries are placed wherever a wide transformation (shuffle)
appears. The DAG Scheduler also applies \emph{partition skipping}: any
stage whose output partitions are already cached in executor storage memory
is skipped entirely, and the cached data is used directly as input to the
next stage.

The \term{Task Scheduler} receives the physical task list for each ready
stage and submits tasks to executors via the cluster manager. The Task
Scheduler applies \emph{delay scheduling}: if the data-local executor is
briefly unavailable, the Task Scheduler waits a configurable period
(default 3 seconds) before falling back to a rack-local or
any-node assignment, maximising data locality.

\begin{figure}[h!t]
\centering
\begin{tikzpicture}[
    box/.style={rectangle, draw=DEBlue!60, fill=DELightBlue!40,
                rounded corners=4pt, minimum width=3.2cm, minimum height=0.75cm,
                font=\small\bfseries\sffamily, align=center},
    cmanager/.style={rectangle, draw=DEAccent!70, fill=DELightAccent,
               rounded corners=4pt, minimum width=3.2cm, minimum height=0.75cm,
               font=\small\bfseries\sffamily, align=center},
    exec/.style={rectangle, draw=DEGreen!60, fill=DELightGreen,
                 rounded corners=4pt, minimum width=2.5cm, minimum height=0.65cm,
                 font=\small\sffamily, align=center},
    task/.style={rectangle, draw=DEGreen!50, fill=DEGreen!10,
                 rounded corners=2pt, minimum width=1.0cm, minimum height=0.45cm,
                 font=\tiny\sffamily, align=center},
    arr/.style={-{Stealth[length=4pt,width=4pt]}, thick},
    lbl/.style={font=\scriptsize\sffamily, align=center}
  ]

  % Driver box (left column)
  \node[box] (ss)  at (0, 7.0) {SparkSession};
  \node[box] (dag) at (0, 5.7) {DAG Scheduler};
  \node[box] (ts)  at (0, 4.4) {Task Scheduler};

  % Cluster Manager (center-right)
  \node[cmanager]  (cm)  at (5.0, 5.7) {Cluster Manager\\(YARN / K8s)};

  % Worker nodes
  \node[exec] (e1) at (3.2, 2.8) {Executor 1};
  \node[exec] (e2) at (5.5, 2.8) {Executor 2};
  \node[exec] (e3) at (7.8, 2.8) {Executor 3};

  % Tasks inside executors
  \node[task] at (2.7, 2.1) {T1}; \node[task] at (3.7, 2.1) {T2};
  \node[task] at (5.0, 2.1) {T3}; \node[task] at (6.0, 2.1) {T4};
  \node[task] at (7.3, 2.1) {T5}; \node[task] at (8.3, 2.1) {T6};

  % Driver internal flow
  \draw[arr, DEBlue] (ss.south)  -- (dag.north)
       node[midway, right, lbl] {\tiny action triggered};
  \draw[arr, DEBlue] (dag.south) -- (ts.north)
       node[midway, right, lbl] {\tiny stage task list};

  % Task Scheduler → Cluster Manager
  \draw[arr, DEAccent] (ts.east) to[out=0, in=-120] (cm.south)
       node[pos=0.5, above, lbl, text=DEAccent] {\tiny request containers};

  % Cluster Manager → Executors
  \draw[arr, DEAccent!70!black] (cm.south) to[out=-90, in=90] (e1.north);
  \draw[arr, DEAccent!70!black] (cm.south) -- (e2.north);
  \draw[arr, DEAccent!70!black] (cm.south) to[out=-90, in=90] (e3.north);
  \node[lbl, text=DEAccent!70!black] at (5.0, 3.9)
       {\tiny launch executor containers};

  % Executors heartbeat to driver
  \draw[arr, DEGray!50, dashed]
       (e2.west) to[bend right=30] (ts.east)
       node[pos=0.4, right, lbl, text=DEGray] {\tiny\itshape status / results};

  % Driver label
  \node[font=\small\bfseries\sffamily, text=DEGray, rotate=90]
       at (-2.2, 5.7) {Driver Process};
  \draw[DEGray!30, dashed, rounded corners]
       (-1.8, 3.9) rectangle (1.8, 7.5);

\end{tikzpicture}
\caption{Spark architecture. The driver process contains the DAG Scheduler
         and Task Scheduler. The cluster manager (YARN, Kubernetes, or Standalone)
         allocates executor containers on worker nodes. Executors run tasks assigned
         by the driver and report results via heartbeat. Executors remain alive for
         the full application lifetime, enabling in-memory data sharing across jobs.}
\label{fig:spark-arch}
\end{figure}

%%====================================================================
\section{Transformations and Actions}
\label{sec:transformations}
%%====================================================================

Every RDD operation in Spark is classified as either a \emph{transformation}
or an \emph{action}. Understanding this distinction, and the further
sub-classification of transformations into narrow and wide, is essential
for writing correct and performant Spark applications.

%--------------------------------------------------------------------
\subsection{Narrow Transformations}
\label{subsec:narrow}
%--------------------------------------------------------------------

A \term{narrow transformation} is one in which each output partition
depends on at most one input partition. No data movement across the
network is required: each executor can compute its output partitions
independently, working only on the input partitions it already holds
locally. Consecutive narrow transformations are \emph{pipelined} by the
DAG Scheduler into a single stage: a task processes a record through all
pipelined transformations without materialising any intermediate RDD.

The most important narrow transformations are:

\begin{itemize}
  \item \texttt{map(f):} applies function $f$ to each record, producing
        one output record per input record.
  \item \texttt{filter(p):} retains only records satisfying predicate $p$.
  \item \texttt{flatMap(f):} applies $f$ to each record and flattens the
        resulting sequences into one output RDD.
  \item \texttt{mapPartitions(f):} applies $f$ once per partition rather
        than once per record, which is more efficient when the function
        has per-call overhead (e.g., opening a database connection or
        loading a model).
  \item \texttt{union(other):} concatenates two RDDs with the same
        schema, with no data movement (each partition of the result is
        exactly one partition of an input).
  \item \texttt{sample(fraction):} samples each partition independently,
        with no cross-partition communication.
\end{itemize}

%--------------------------------------------------------------------
\subsection{Wide Transformations and the Shuffle}
\label{subsec:wide}
%--------------------------------------------------------------------

A \term{wide transformation} is one in which each output partition may
depend on multiple input partitions, potentially held on different
executors. Producing the output requires \emph{moving data across the
network}---a \term{shuffle}. Every shuffle is a stage boundary: the
DAG Scheduler inserts a synchronisation point before the downstream
stage begins, ensuring all upstream tasks have completed and written
their shuffle output before the downstream tasks begin reading it.

The shuffle is the most expensive operation in Spark. Each shuffle
involves:
\begin{enumerate}
  \item \textbf{Map side (shuffle write):} each upstream task
        hash-partitions its output records by key and writes each
        partition's records to a local shuffle file on the executor's
        disk. This write is unavoidable: the data must survive until
        all downstream tasks have read it.
  \item \textbf{Network transfer:} each downstream task reads the shuffle
        blocks it needs from remote executors via HTTP, crossing the
        network.
  \item \textbf{Reduce side (shuffle read):} each downstream task
        receives and potentially sorts its input before applying the
        aggregation or join function.
\end{enumerate}

The most important wide transformations are:

\begin{itemize}
  \item \texttt{groupByKey():} collects all values for each key into a
        single iterable in one partition. Requires a full shuffle.
  \item \texttt{reduceByKey(f):} applies the commutative and associative
        function $f$ to reduce values for each key. Crucially,
        \texttt{reduceByKey} applies a partial aggregation \emph{within
        each partition} before the shuffle (analogous to the Combiner in
        MapReduce), dramatically reducing the volume of data transferred.
  \item \texttt{aggregateByKey(zeroValue)(seqOp, combOp):} a generalisation
        of \texttt{reduceByKey} that supports different input and output types.
  \item \texttt{sortByKey(ascending):} globally sorts the RDD by key,
        requiring a range-partitioned shuffle.
  \item \texttt{join(other):} joins two RDDs on their keys, requiring
        both to be hash-shuffled by key to co-locate matching records.
  \item \texttt{distinct():} removes duplicates, requiring a shuffle to
        co-locate all copies of each record.
  \item \texttt{repartition(n):} redistributes the RDD into exactly $n$
        partitions, performing a full shuffle.
\end{itemize}

\begin{figure}[h!t]
\centering
\begin{tikzpicture}[
    part/.style={rectangle, draw=DEBlue!50, fill=DELightBlue!50,
                 minimum width=1.6cm, minimum height=0.6cm,
                 font=\tiny\sffamily, align=center},
    shuffle/.style={rectangle, draw=DEAccent!70, fill=DELightAccent,
                    minimum width=1.6cm, minimum height=0.6cm,
                    font=\tiny\sffamily, align=center},
    stage/.style={font=\small\bfseries\sffamily, text=DEBlue},
    arr/.style={-{Stealth[length=3.5pt,width=3.5pt]}, thick},
    lbl/.style={font=\scriptsize\sffamily, align=center}
  ]

  % Stage 1
  \node[stage] at (1.2, 6.2) {Stage 1};
  \node[part] (s1p1) at (0,   5.5) {textFile\\$P_1$};
  \node[part] (s1p2) at (1.6, 5.5) {textFile\\$P_2$};
  \node[part] (s1p3) at (3.2, 5.5) {textFile\\$P_3$};

  % Narrow transforms (pipelined)
  \node[lbl, text=DEGreen!70!black] at (-1.5, 4.5)
       {flatMap, map\\(narrow, pipelined)};
  \draw[DEGreen!60, dashed] (-0.8, 4.3) -- (0, 4.8);

  \node[part] (s1o1) at (0,   4.5) {map out\\$P_1$};
  \node[part] (s1o2) at (1.6, 4.5) {map out\\$P_2$};
  \node[part] (s1o3) at (3.2, 4.5) {map out\\$P_3$};

  \draw[arr, DEBlue!60] (s1p1.south) -- (s1o1.north);
  \draw[arr, DEBlue!60] (s1p2.south) -- (s1o2.north);
  \draw[arr, DEBlue!60] (s1p3.south) -- (s1o3.north);

  % Shuffle boundary
  \draw[DERed!60, very thick, dashed]
       (-0.9, 3.9) -- (4.1, 3.9);
  \node[lbl, text=DERed] at (5.4, 3.9)
       {\textbf{Shuffle boundary}\\(reduceByKey)};

  % Stage 2
  \node[stage] at (1.2, 3.4) {Stage 2};

  \node[shuffle] (s2p1) at (0,   3.1) {reduce\\$K_1$};
  \node[shuffle] (s2p2) at (1.6, 3.1) {reduce\\$K_2$};
  \node[shuffle] (s2p3) at (3.2, 3.1) {reduce\\$K_3$};

  % All-to-all shuffle arrows
  \foreach \src in {s1o1, s1o2, s1o3}
    \foreach \dst in {s2p1, s2p2, s2p3}
      \draw[DEAccent!50, thin] (\src.south) -- (\dst.north);

  % Output
  \node[part] (out1) at (0,   2.1) {out $P_1$};
  \node[part] (out2) at (1.6, 2.1) {out $P_2$};
  \node[part] (out3) at (3.2, 2.1) {out $P_3$};
  \draw[arr, DEGreen!70!black] (s2p1.south) -- (out1.north);
  \draw[arr, DEGreen!70!black] (s2p2.south) -- (out2.north);
  \draw[arr, DEGreen!70!black] (s2p3.south) -- (out3.north);
  \node[lbl, text=DEGreen!70!black] at (-1.5, 2.2) {saveAsTextFile\\(action)};
  \draw[DEGreen!60, dashed] (-0.8, 2.2) -- (0, 2.2);

\end{tikzpicture}
\caption{Spark stage and shuffle boundary for a word-count job. Stage~1 pipelines
         \texttt{textFile}, \texttt{flatMap}, and \texttt{map} (all narrow) into a
         single pass. The \texttt{reduceByKey} is a wide transformation, creating a
         shuffle boundary. Stage~2 performs the shuffle read and final aggregation.
         The all-to-all arrows represent network transfer of shuffle blocks.}
\label{fig:shuffle-boundary}
\end{figure}

%--------------------------------------------------------------------
\subsection{groupByKey vs.\ reduceByKey: A Critical Distinction}
\label{subsec:groupby-vs-reduceby}
%--------------------------------------------------------------------

\begin{caution}{\texttt{groupByKey} Transfers All Data; \texttt{reduceByKey} Does Not}
The most common performance anti-pattern in Spark is using
\texttt{groupByKey()} when \texttt{reduceByKey()} is applicable.
Consider computing the sum of values per key. \texttt{groupByKey()}
sends \emph{all} raw values for each key across the network to one
executor, where the sum is computed. If there are $10^8$ records with
key ``DE'' each holding a 16-byte value, that is 1.6\GB\ of network
traffic \emph{just for that one key}.

\texttt{reduceByKey(lambda a, b: a + b)} first sums values within each
source partition independently (producing one sum per partition per key),
then shuffles only these partial sums. If there are 200 partitions,
the shuffle transfers 200 numbers for key ``DE''---16-byte numbers---
instead of 1.6\GB. The network amplification factor is reduced by the
partition count ($200\times$ in this example).

The general rule: use \texttt{reduceByKey}, \texttt{aggregateByKey}, or
\texttt{combineByKey} whenever the aggregation function is associative
and commutative. Use \texttt{groupByKey} only when all raw values for a
key must be available simultaneously (e.g., computing a per-key median).
\end{caution}

%--------------------------------------------------------------------
\subsection{Actions}
\label{subsec:actions}
%--------------------------------------------------------------------

\term{Actions} are the operations that trigger the execution of the
accumulated RDD lineage DAG and return a concrete result. The most
important actions are:

\begin{itemize}
  \item \texttt{count():} returns the number of records in the RDD as a
        \texttt{Long}.
  \item \texttt{collect():} transfers all partitions to the driver and
        returns a local array (use only for small results).
  \item \texttt{take(n):} returns the first $n$ records from the first
        partition(s), with less data movement than \texttt{collect()}.
  \item \texttt{reduce(f):} applies the associative and commutative
        function $f$ to all records, returning a single value.
  \item \texttt{foreach(f):} applies function $f$ to each record on
        the executor side (e.g., to write records to an external system)
        without returning data to the driver.
  \item \texttt{saveAsTextFile(path):} writes the RDD as plain text
        files to an HDFS path (one file per partition).
  \item \texttt{saveAsSequenceFile(path):} writes as Hadoop SequenceFile
        format (key-value pairs).
\end{itemize}

%%====================================================================
\section{Spark SQL: DataFrames, Datasets, and the Catalyst Optimizer}
\label{sec:spark-sql}
%%====================================================================

The RDD API provides low-level control over distributed computation but
requires the programmer to reason explicitly about types, partitions, and
data layouts. For SQL-style analytical workloads, this level of control
is unnecessarily cumbersome. Spark SQL, introduced in Spark 1.0 and
substantially redesigned through Spark 2.0 and 3.0, provides a
structured, schema-aware API that allows the Spark engine to apply
far more aggressive automatic optimisations.

%--------------------------------------------------------------------
\subsection{DataFrames and Datasets}
\label{subsec:dataframes}
%--------------------------------------------------------------------

A \term{DataFrame} is an RDD of \texttt{Row} objects with a named,
typed schema---conceptually a distributed relational table. The schema
is declared explicitly or inferred from structured data sources (Parquet,
ORC, JSON, Hive tables, JDBC). Unlike a raw RDD, a DataFrame carries
column names and data types, enabling the query engine to reason about
the computation without executing it.

A \term{Dataset} is a strongly typed, JVM-object-centric version of
the DataFrame, available in Scala and Java. A \texttt{Dataset[Person]}
stores JVM \texttt{Person} objects rather than generic \texttt{Row}
objects, providing compile-time type safety at the cost of some JVM
serialisation overhead. In Python and R, DataFrames are used exclusively
(there is no separate Dataset API because these languages are dynamically
typed).

The \term{SparkSession} is the unified entry point for all Spark SQL
operations:

\begin{lstlisting}[style=pyspark, caption={SparkSession creation and DataFrame operations}]
from pyspark.sql import SparkSession
from pyspark.sql.functions import col, avg, count

spark = SparkSession.builder \
    .appName("SalesAnalysis") \
    .config("spark.sql.shuffle.partitions", "200") \
    .getOrCreate()

# Load a Parquet table from HDFS
sales = spark.read.parquet("hdfs:///data/sales/dt=2024-*")

# DataFrame operations: filter, groupBy, aggregate
result = (sales
    .filter(col("amount") > 100)
    .groupBy("region", "product_id")
    .agg(
        count("*").alias("num_sales"),
        avg("amount").alias("avg_amount")
    )
    .orderBy(col("avg_amount").desc())
)

result.write.mode("overwrite").parquet("hdfs:///output/sales_summary")
\end{lstlisting}

%--------------------------------------------------------------------
\subsection{The Catalyst Optimizer}
\label{subsec:catalyst}
%--------------------------------------------------------------------

The key advantage of the DataFrame/SQL API over the raw RDD API is that
every DataFrame operation passes through the \term{Catalyst optimizer}
before execution. Catalyst is a tree-based query optimizer written in
Scala that applies a cascade of rule-based and cost-based transformations
to reduce query execution cost. It operates in four phases.

\textbf{Phase 1 --- Analysis.} Catalyst resolves attribute names and
data types by consulting the \term{Spark Catalog}---the metadata registry
for tables, views, functions, and databases (backed by the Hive Metastore
in production). Unresolved logical plans (trees with unresolved attribute
references) become resolved logical plans (trees where every attribute
has a confirmed name and data type).

\textbf{Phase 2 --- Logical Optimisation.} Catalyst applies rule-based
algebraic transformations to the resolved logical plan:
\begin{itemize}
  \item \textbf{Predicate pushdown:} move \texttt{WHERE} filter conditions
        as close to the data source as possible, so that Parquet row
        groups can be skipped before data is read into memory.
  \item \textbf{Column pruning:} remove columns not referenced by the
        query from the scan; only referenced column chunks are read
        from Parquet files.
  \item \textbf{Constant folding:} evaluate expressions that involve
        only constants at compile time (e.g., \texttt{WHERE year = 2020
        + 4} becomes \texttt{WHERE year = 2024}).
  \item \textbf{Join reordering:} place smaller tables on the build side
        of hash joins (or as the broadcast side of broadcast hash joins).
\end{itemize}

\textbf{Phase 3 --- Physical Planning.} Catalyst converts the optimised
logical plan into one or more candidate physical plans---concrete
execution strategies involving Spark operators. For a join, physical
plans may include sort-merge join (for large tables), broadcast hash join
(when one table fits in executor memory, default threshold 10\MB), or
hash join. Catalyst applies a simple cost model (based on estimated row
counts from Parquet file statistics) to select the lowest-cost plan.

\textbf{Phase 4 --- Code Generation (Tungsten).} The selected physical
plan is compiled to JVM bytecode by the \term{Tungsten} execution engine.
Tungsten uses \emph{whole-stage code generation}: rather than interpreting
a query plan row-by-row through a generic operator tree, Tungsten
generates a single JVM method that fuses all operators in a pipeline into
a tight loop. This eliminates virtual function call overhead and enables
the JIT compiler to produce CPU-efficient native code. Tungsten also uses
off-heap \emph{unsafe} memory for sort and aggregate buffers, bypassing
JVM garbage collection overhead.

\begin{systemdesign}{Catalyst and Parquet: An End-to-End Example}
A query \texttt{SELECT product\_id, SUM(amount) FROM sales WHERE region
= 'EU' AND dt BETWEEN '2024-01-01' AND '2024-03-31' GROUP BY
product\_id} on a Parquet table partitioned by \texttt{dt} passes
through Catalyst as follows:
\begin{enumerate}
  \item \textbf{Partition pruning:} only the subdirectories
        \texttt{/sales/dt=2024-01-01/} through
        \texttt{/sales/dt=2024-03-31/} are scanned. All other partitions
        are eliminated before any data is read.
  \item \textbf{Predicate pushdown:} the \texttt{region = 'EU'} filter
        is pushed into the Parquet reader. Row groups where
        \texttt{max(region) < 'EU'} or \texttt{min(region) > 'EU'}
        (lexicographically) are skipped via footer statistics.
  \item \textbf{Column pruning:} only the \texttt{region},
        \texttt{product\_id}, and \texttt{amount} column chunks are read;
        all other column chunks are not opened.
  \item \textbf{Aggregation:} Catalyst generates a two-phase aggregate
        plan: a partial sum within each partition (no shuffle), then a
        shuffle on \texttt{product\_id}, then a final sum per partition
        in the reduce stage.
\end{enumerate}
The combined effect of partition pruning, predicate pushdown, and column
pruning can reduce the physical I/O from terabytes to gigabytes on a
selective query---without any manual hint or configuration beyond
the initial table partitioning.
\end{systemdesign}

%--------------------------------------------------------------------
\subsection{SparkSQL and Hive Integration}
\label{subsec:spark-hive}
%--------------------------------------------------------------------

Spark SQL integrates with the \tool{Hive Metastore} through the Hive
Catalog interface, allowing SparkSQL queries to operate on tables defined
by Hive DDL statements, including partitioned Parquet and ORC tables
(introduced in Chapter~\ref{ch:yarn-formats}). This integration allows
organisations to migrate from Hive-on-MapReduce to Spark SQL with no
changes to table definitions, partition schemes, or ETL job logic---only
the execution engine changes.

%%====================================================================
\section{MLlib: Distributed Machine Learning}
\label{sec:mllib}
%%====================================================================

Apache Spark MLlib is Spark's built-in machine learning library. Its
central design goal is to express machine learning pipelines---sequences
of data preprocessing, feature engineering, and model training steps---as
\term{Pipeline} objects that can be fit to training data, saved to disk,
and applied to new data with a single API call.

%--------------------------------------------------------------------
\subsection{Core Abstractions: Transformer, Estimator, and Pipeline}
\label{subsec:pipeline-api}
%--------------------------------------------------------------------

The MLlib Pipeline API is built from three abstractions.

A \term{Transformer} is an object that implements a \texttt{transform(df)}
method: it takes a DataFrame as input and produces a new DataFrame with
one or more new columns added. Transformers are \emph{stateless}---they
do not learn parameters from data. Examples include \texttt{Tokenizer}
(splits text into words), \texttt{HashingTF} (converts token lists to
TF-IDF feature vectors), \texttt{StandardScaler} after fitting (normalises
numerical features using pre-computed mean and standard deviation), and
any trained model (which transforms a feature DataFrame into a
predictions DataFrame).

An \term{Estimator} is an object that implements a \texttt{fit(df)}
method: it trains a model on data and produces a \emph{Transformer}
(the trained model). Examples include \texttt{LogisticRegression},
\texttt{RandomForestClassifier}, \texttt{KMeans}, and
\texttt{StandardScaler} (before fitting---the fitted version is a
Transformer).

A \term{Pipeline} is an ordered sequence of Transformers and Estimators.
\texttt{pipeline.fit(df)} calls \texttt{fit()} on each Estimator in
sequence, using the training DataFrame and passing the output to the
next stage. The result is a \texttt{PipelineModel}---a sequence of
Transformers---which can be applied to new data via
\texttt{model.transform(test\_df)}.

\begin{lstlisting}[style=pyspark, caption={MLlib Pipeline for text classification}]
from pyspark.ml import Pipeline
from pyspark.ml.feature import (Tokenizer, HashingTF, IDF,
                                 StringIndexer, VectorAssembler)
from pyspark.ml.classification import LogisticRegression
from pyspark.ml.evaluation import MulticlassClassificationEvaluator

# Load labelled text data
df = spark.read.parquet("hdfs:///data/news_articles")

# Build pipeline stages
tokenizer  = Tokenizer(inputCol="text", outputCol="tokens")
hashing_tf = HashingTF(inputCol="tokens", outputCol="raw_features",
                       numFeatures=20_000)
idf        = IDF(inputCol="raw_features", outputCol="features")
indexer    = StringIndexer(inputCol="label", outputCol="label_idx")
lr         = LogisticRegression(featuresCol="features",
                                labelCol="label_idx",
                                maxIter=100)

pipeline = Pipeline(stages=[tokenizer, hashing_tf, idf, indexer, lr])

# Train / test split
train, test = df.randomSplit([0.8, 0.2], seed=42)

# Fit: trains IDF, StringIndexer, LogisticRegression on train
model = pipeline.fit(train)

# Transform: applies all transformations to test
predictions = model.transform(test)

evaluator = MulticlassClassificationEvaluator(
    labelCol="label_idx", predictionCol="prediction",
    metricName="accuracy")
print(f"Test accuracy: {evaluator.evaluate(predictions):.4f}")

# Persist the trained model to HDFS for serving
model.save("hdfs:///models/news_classifier_v1")
\end{lstlisting}

%--------------------------------------------------------------------
\subsection{Key MLlib Algorithms}
\label{subsec:mllib-algos}
%--------------------------------------------------------------------

\begin{table}[h!t]
\centering
\caption{Representative MLlib algorithms by category}
\label{tab:mllib-algos}
\renewcommand{\arraystretch}{1.5}
\begin{tabular}{L{2.4cm} L{3.6cm} L{4.4cm}}
\toprule
\textbf{Category} & \textbf{Algorithm} & \textbf{Typical use case} \\
\midrule
Classification  & LogisticRegression, RandomForestClassifier, GBTClassifier, LinearSVC & Spam detection, churn prediction, document classification \\
Regression      & LinearRegression, RandomForestRegressor, GBTRegressor & Price forecasting, demand estimation \\
Clustering      & KMeans, GaussianMixture, BisectingKMeans & Customer segmentation, anomaly detection \\
Recommendation  & ALS (Alternating Least Squares) & Collaborative filtering (Netflix, Spotify) \\
Feature Eng.    & PCA, Word2Vec, MinHashLSH, QuantileDiscretizer & Dimensionality reduction, similarity search \\
\bottomrule
\end{tabular}
\end{table}

\textbf{Alternating Least Squares (ALS)} deserves particular mention as
it illustrates how Spark's in-memory model enables iterative ML algorithms.
ALS is a matrix factorisation algorithm for collaborative filtering: given
a user-item rating matrix with many missing entries, it alternately holds
user factors fixed and solves for item factors (an embarrassingly parallel
linear algebra problem), then holds item factors fixed and solves for user
factors. Each alternation reads the full ratings matrix. On a dataset of
$10^9$ ratings, this read is performed 10--20 times per training run; with
the ratings cached in Spark's distributed memory, all subsequent reads are
served from RAM at disk-free speed.

%%====================================================================
\section{Memory Management in Spark}
\label{sec:memory}
%%====================================================================

Spark's performance advantage over MapReduce rests on its ability to
use the executor's JVM heap for both active computation and RDD caching.
Understanding how Spark divides this memory between competing uses is
essential for diagnosing and resolving the two most common production
failure modes: spill-to-disk (which degrades performance) and OOM
(OutOfMemory) errors (which crash executors).

%--------------------------------------------------------------------
\subsection{The Unified Memory Model}
\label{subsec:unified-mem}
%--------------------------------------------------------------------

Spark 1.6 introduced the \term{unified memory model}, which replaced
the earlier static partitioning of execution and storage memory with a
dynamic boundary. The executor JVM heap is divided into three regions:

\begin{enumerate}
  \item \textbf{Reserved Memory} (300\MB, fixed): reserved for Spark's
        internal objects. Not configurable.
  \item \textbf{Spark Memory} (\texttt{spark.memory.fraction} $\times$
        (heap $-$ 300\MB), default $0.6$): the pool shared between
        execution and storage.
  \item \textbf{User Memory} (remaining $1 - 0.6 = 0.4$ fraction):
        available for user-defined data structures, UDF objects, and
        the Java/Scala collections used in application code.
\end{enumerate}

Within the Spark Memory pool, execution memory and storage memory share
the space dynamically:

\begin{itemize}
  \item \textbf{Storage Memory} (initially \texttt{spark.memory.storageFraction}
        $\times$ Spark Memory, default $0.5$) caches persisted RDD
        partitions. Storage memory can expand into unused execution memory,
        but cannot forcibly evict active execution memory.
  \item \textbf{Execution Memory} (the remainder) holds sort buffers,
        hash tables for aggregations, and shuffle read/write buffers.
        Execution memory can borrow from storage memory, and if it does,
        Spark evicts cached RDD partitions (via LRU) to make room.
\end{itemize}

\begin{defbox}{Unified Memory Allocation}
For an executor with $H$ bytes of JVM heap:
\begin{align}
M_{\text{spark}}   &= \texttt{fraction} \times (H - 300\MB) \label{eq:spark-mem} \\
M_{\text{storage}} &\approx \texttt{storageFraction} \times M_{\text{spark}} \label{eq:storage-mem} \\
M_{\text{exec}}    &= M_{\text{spark}} - M_{\text{storage}} \quad \text{(initial)} \label{eq:exec-mem}
\end{align}
For $H = 16\GB$, \texttt{fraction} $= 0.6$, \texttt{storageFraction}
$= 0.5$:
\begin{align*}
M_{\text{spark}}   &= 0.6 \times (16\GB - 0.3\GB) \approx 9.4\GB \\
M_{\text{storage}} &\approx 0.5 \times 9.4\GB = 4.7\GB \\
M_{\text{exec}}    &\approx 4.7\GB \quad \text{(can grow up to } 9.4\GB)
\end{align*}
\end{defbox}

%--------------------------------------------------------------------
\subsection{Serialisation and Kryo}
\label{subsec:kryo}
%--------------------------------------------------------------------

When RDD partitions are persisted with \texttt{MEMORY\_ONLY\_SER} or
spill to disk, Spark must serialise Java/Scala objects into a byte
stream. By default, Spark uses Java's native serialisation, which is
flexible (handles arbitrary Java objects) but slow and space-inefficient.

\term{Kryo serialisation} (\texttt{spark.serializer = org.apache.spark.
serializer.KryoSerializer}) is a third-party library that serialises
registered classes 2--10$\times$ faster and produces 3--5$\times$ smaller
byte representations than Java serialisation. For RDDs of simple, known
types (such as \texttt{(String, Int)} pairs or custom case classes),
Kryo is almost always the correct choice.

\begin{lstlisting}[style=pyspark, caption={Enabling Kryo serialization and registering classes}]
from pyspark import SparkConf, SparkContext

conf = SparkConf() \
    .setAppName("KryoDemo") \
    .set("spark.serializer",
         "org.apache.spark.serializer.KryoSerializer") \
    .set("spark.kryo.registrationRequired", "false") \
    .set("spark.kryoserializer.buffer.max", "256m")

sc = SparkContext(conf=conf)
\end{lstlisting}

%--------------------------------------------------------------------
\subsection{Data Skew and Salting}
\label{subsec:skew}
%--------------------------------------------------------------------

\term{Data skew} occurs when the keys of a wide transformation are
distributed unevenly: a small number of keys contain a very large
fraction of the records. In a \texttt{reduceByKey}, all records for
the same key must be sent to the same partition. If 40\% of the dataset
has the key \texttt{``US''}, one task will process 40\% of the data
while the other 199 tasks process the remaining 60\% divided evenly.
The entire stage cannot complete until the skewed task finishes. The
stage latency is dominated by the slowest task, making the overall
job wall-clock time proportional to the largest single-key volume.

The standard remedy is \term{salting}: appending a random integer suffix
(the ``salt'') to skewed keys to distribute their records across multiple
partitions, then performing a two-phase aggregation.

\begin{lstlisting}[style=pyspark, caption={Salting to resolve data skew in a PySpark RDD}]
import random

SALT_FACTOR = 50    # number of salt buckets for skewed keys
SKEWED_KEYS = {"US", "CN", "DE"}   # keys known to be skewed

def add_salt(record):
    key, value = record
    if key in SKEWED_KEYS:
        salted_key = (key, random.randint(0, SALT_FACTOR - 1))
    else:
        salted_key = (key, 0)
    return (salted_key, value)

def remove_salt(record):
    (key, _salt), value = record
    return (key, value)

# Phase 1: shuffle by (key, salt) -> 50x more reducers for skewed keys
partial = (rdd
    .map(add_salt)
    .reduceByKey(lambda a, b: a + b))   # partial aggregation

# Phase 2: strip salt, re-aggregate to get final per-key totals
final = (partial
    .map(remove_salt)
    .reduceByKey(lambda a, b: a + b))   # final aggregation
\end{lstlisting}

\begin{techdebt}{Skew is Invisible to the Optimizer}
Data skew is one of the most common causes of Spark job failure in
production, and it is almost entirely invisible to the Catalyst optimizer.
Catalyst can pushdown predicates, prune columns, and choose between sort-merge
and broadcast joins---but it cannot detect that 40\% of the data has the
same key unless the user provides statistics or explicit hints. The symptoms
are: a stage where 199 out of 200 tasks complete quickly but one task runs
for hours, eventually causing the executor to OOM or timeout.

Diagnostic steps: examine the Spark UI ``Stages'' tab and look for tasks with
anomalously large ``Input Size'' or ``Shuffle Read Size.'' Use
\texttt{df.groupBy("key").count().orderBy(F.desc("count")).show(20)} to identify
skewed keys before designing the aggregation.
\end{techdebt}

%%====================================================================
\section{Research Spotlights}
\label{sec:ch6-research}
%%====================================================================

\begin{researchspotlight}{Zaharia et al.\ (2010) --- Spark: Cluster Computing with Working Sets}
\textbf{Citation:} Zaharia, M., Chowdhury, M., Franklin, M.~J.,
Shenker, S., and Stoica, I.\ (2010). Spark: Cluster Computing with
Working Sets. \textit{Proceedings of the 2nd USENIX Workshop on Hot
Topics in Cloud Computing (HotCloud 2010)}.
\parencite{Spark2010}

\smallskip\noindent
\textbf{Key insight:} The paper's thesis is that MapReduce is fundamentally
unsuited for two important and growing workload classes: \emph{iterative
algorithms} (machine learning, graph processing) and \emph{interactive
data analysis} (exploratory SQL on cached datasets). Both require
reading the same data multiple times, and MapReduce's forced disk
materialisation between stages makes this prohibitively expensive.
The solution is to treat distributed RAM as a first-class abstraction
that persists across job boundaries.

\smallskip\noindent
\textbf{Technical contributions:}
\begin{itemize}
  \item Introduced the concept of a \emph{working set}---the subset of
        data that an application accesses repeatedly---and argued that
        a cluster computing framework should allow the working set to be
        stored in distributed memory across the cluster.
  \item Demonstrated that a logistic regression training job ran
        $10\times$ faster on Spark than on Hadoop MapReduce by caching
        the training set in memory after the first iteration.
  \item Showed that interactive queries over a 39\GB\ Wikipedia dataset
        ran in 0.5--1 second on cached data vs.\ 20--60 seconds on
        Hadoop, enabling genuine interactive exploration at scale.
\end{itemize}

\smallskip\noindent
\textbf{Impact today:} The HotCloud 2010 paper introduced Spark to the
research community and has been cited over 5{,}000 times. Within three
years of this publication, Spark had been adopted by hundreds of
organisations and had become the fastest-growing Apache project in
history. The working-set insight---that most production ML and analytics
workloads repeatedly access the same data---is now the primary
justification for in-memory computing frameworks across the industry.
\end{researchspotlight}

\begin{researchspotlight}{Zaharia et al.\ (2012) --- Resilient Distributed Datasets: A Fault-Tolerant Abstraction for In-Memory Cluster Computing}
\textbf{Citation:} Zaharia, M., Chowdhury, M., Das, T., Dave, A., Ma, J.,
McCauley, M., Franklin, M.~J., Shenker, S., and Stoica, I.\ (2012).
Resilient Distributed Datasets: A Fault-Tolerant Abstraction for In-Memory
Cluster Computing. \textit{Proceedings of the 9th USENIX Symposium on
Networked Systems Design and Implementation (NSDI 2012)}, pp.\ 15--28.
\textbf{Best Paper Award.} \parencite{RDD2012}

\smallskip\noindent
\textbf{Key insight:} The central contribution is the formalization of
the RDD abstraction and its lineage-based fault tolerance model. The
paper argues that existing distributed shared memory systems fail for
large-scale data processing because they require fine-grained writes to
be tracked for replication---an overhead that is prohibitive at
petabyte scale. RDDs restrict mutations to coarse-grained
transformations (each transformation applies the same function to all
records in a partition), which makes it possible to describe any lost
partition's provenance with a short lineage graph rather than a full
data log.

\smallskip\noindent
\textbf{Technical contributions:}
\begin{itemize}
  \item Formally defined the RDD as a five-tuple (partitions, dependencies,
        compute function, partitioner, preferred locations) and proved
        that lineage-based fault tolerance is sufficient for the class
        of coarse-grained transformations that covers all common data
        analytics operations.
  \item Distinguished narrow from wide dependencies and showed that
        narrow-dependency lineage recomputation is always partition-local
        (no cross-executor coordination needed), while wide-dependency
        recomputation may require reading multiple parent partitions.
  \item Benchmarked Spark RDDs against Hadoop MapReduce on four workloads:
        logistic regression (25$\times$ speedup on second iteration),
        PageRank (7$\times$ speedup), $k$-means clustering (40$\times$
        speedup), and interactive SQL queries (sub-second vs.\ tens of
        seconds).
\end{itemize}

\smallskip\noindent
\textbf{Impact today:} The NSDI 2012 paper is the most-cited Spark
paper, with over 10{,}000 citations, and is considered one of the
foundational papers of the systems era of big data processing. The RDD
abstraction directly influenced the design of DataFrame/Dataset APIs,
Apache Flink's DataSet/DataStream APIs, TensorFlow's distributed variable
model, and Ray's distributed object store. The narrow/wide dependency
distinction is now standard vocabulary in distributed systems courses
worldwide.
\end{researchspotlight}

%%====================================================================
\section{Case Study: Alibaba's Spark Deployment for Singles' Day (11.11)}
\label{sec:alibaba-case}
%%====================================================================

\begin{casestudy}{Alibaba --- Apache Spark at 10{,}000 Nodes for the World's Largest Shopping Event}

\textbf{Background.}
Alibaba's Singles' Day (11.11) shopping festival---held annually on
11 November---is the world's largest single-day commerce event.
In 2019, the event generated \$38.4 billion USD in gross merchandise
value within 24 hours, with a peak transaction rate of 544{,}000
orders per second. Behind this scale lay one of the largest production
deployments of Apache Spark in existence, integrated into Alibaba's
MaxCompute (formerly ODPS) and E-MapReduce platforms.

\smallskip
\textbf{Scale.}
Alibaba's internal Spark deployment for 11.11 2019 involved:
\begin{itemize}
  \item Over \textbf{10{,}000 executor nodes}, each with 96\GB\ of RAM
        and 24 vCPUs---approximately 960\TB\ of aggregate RAM and
        240{,}000 vCPUs available to Spark applications.
  \item Processing over \textbf{970\PB} of data in the 24-hour event
        window, primarily through Spark SQL queries over Parquet tables
        partitioned by buyer\_id, seller\_id, and event timestamp.
  \item Over \textbf{1 trillion individual user behaviour events}
        (page views, searches, cart additions, purchases, refunds)
        accumulated, processed, and fed into real-time recommendation
        engines during the event.
\end{itemize}

\smallskip
\textbf{The recommendation challenge.}
The highest-impact use of Spark on 11.11 was powering Taobao's
product recommendation system, which generates personalised product
lists for 800 million active users. Recommendation model updates during
the event must incorporate behaviour signals from the event itself:
a user who viewed three laptops in the first hour of the event should
be shown accessories in the second hour, not laptops she has already
browsed.

This requirement---updating models on streaming data \emph{during} the
event, not just in nightly batch jobs---drove Alibaba's engineers to
build a hybrid streaming-batch architecture:

\begin{enumerate}
  \item \textbf{Behaviour ingestion} (Alibaba's Blink / Apache Flink):
        raw user events arrive from Taobao at approximately 11 million
        events per second at peak. Flink processes them in real time,
        computing streaming feature updates (user-product affinity scores,
        click-through rates, conversion funnels) with a latency of
        under 500 milliseconds.

  \item \textbf{Model refresh} (Spark MLlib): every 10 minutes during
        the 24-hour event, a Spark MLlib ALS collaborative filtering job
        is triggered, consuming the latest streaming feature snapshots
        from HDFS (written by Flink) as training data. The 10-minute
        ALS job runs on 2{,}000 executor nodes with the rating matrix
        (approximately 800 million users $\times$ 2 billion items, with
        40 billion observed entries) cached in distributed memory across
        the executors. Retraining from scratch on this matrix would
        require 6--8 hours with disk-based MapReduce; with Spark and
        in-memory caching, the job completes in 8--12 minutes.

  \item \textbf{Recommendation serving}: the refreshed model parameters
        (user and item latent factors, each a 128-dimensional vector)
        are pushed to Alibaba's distributed key-value store and served
        to Taobao's API layer within 2 minutes of the ALS job completing.
        Users browsing during the refresh window see recommendations based
        on model parameters that incorporate signals from as recently as
        12 minutes ago.
\end{enumerate}

\smallskip
\textbf{Spark SQL for campaign analytics.}
Beyond recommendation, Alibaba's marketing analytics team ran Spark SQL
queries continuously during 11.11 to power the event's live broadcast
analytics dashboard---visible to 800 million TV viewers as a real-time
counter of GMV, transaction volume, and category breakdowns. These
queries joined three Parquet tables:

\begin{center}
\begin{tabular}{lll}
\toprule
\textbf{Table} & \textbf{Size} & \textbf{Update frequency} \\
\midrule
\texttt{transactions} & $\sim$200\TB\ at event end & Appended per minute \\
\texttt{products}     & $\sim$3\GB\ (broadcast)    & Static \\
\texttt{sellers}      & $\sim$5\GB\ (broadcast)    & Static \\
\bottomrule
\end{tabular}
\end{center}

Catalyst automatically applied \emph{broadcast hash join} for the
small dimension tables (products, sellers), avoiding shuffles entirely
for those joins. The \texttt{transactions} table, partitioned by
\texttt{event\_minute}, was scanned with partition pruning: each
dashboard query only read the most recent 5-minute partition
($\sim$600\MB) rather than the full 200\TB\ of accumulated transactions.

\smallskip
\textbf{Skew management.}
A critical engineering challenge was the Taobao flagship stores: a
single store operated by a major brand might account for hundreds of
millions of transaction records in 24 hours. Grouping by
\texttt{seller\_id} without salting would have created severe skew,
with one task processing $>$20\% of the entire dataset. Alibaba's
engineering team applied salting for known high-volume sellers, reducing
the maximum single-task volume to approximately 2\% of the dataset.

\smallskip
\textbf{Outcomes and lessons.}
Alibaba's 11.11 2019 Spark deployment demonstrated three principles
that have since become industry standards for large-scale event analytics:

\begin{itemize}
  \item \textbf{In-memory iterative retraining is viable at billion-user
        scale.} The ALS rating matrix caching strategy reduced model
        refresh time by $40\times$ compared to disk-based alternatives,
        enabling 144 model refreshes during the 24-hour event.
  \item \textbf{Broadcast joins for dimension tables eliminate shuffle.}
        Alibaba measured a $15\times$ latency reduction for multi-table
        queries by broadcasting the 3\GB\ products table to all 10{,}000
        executor nodes, compared to sort-merge joining it through a
        shuffle.
  \item \textbf{Partition pruning must be designed into the schema.}
        Partitioning the transactions table by \texttt{event\_minute}
        rather than \texttt{dt} (which is standard for daily analytics)
        allowed dashboard queries to scan only the most recent partition.
        This was a deliberate schema design decision made 6 weeks before
        the event; retroactively repartitioning 200\TB\ during the event
        would have been impossible.
\end{itemize}
\end{casestudy}

%%====================================================================
\section*{Chapter Summary}
\label{sec:ch6-summary}
%%====================================================================

\begin{itemize}
  \item MapReduce's disk I/O bottleneck for iterative algorithms is
        quantified by the formula: total I/O $= I \times D \times r$,
        where $I$ is the number of iterations, $D$ is the dataset size,
        and $r$ is the HDFS replication factor. For 100 iterations
        on 100\GB\ with $r=3$, this produces 30\TB\ of disk I/O.

  \item An \textbf{RDD} is an immutable, partitioned, lazily computed,
        fault-tolerant collection of records. Its five defining properties
        are: partitions, dependencies (lineage), compute function,
        partitioning metadata, and preferred locations.

  \item \textbf{Lineage-based fault tolerance} recomputes lost partitions
        from their lineage graph rather than maintaining replicas. This
        costs nothing in memory but may require recomputation after
        long lineage chains; \textbf{checkpointing} truncates the lineage
        to bound recomputation cost.

  \item \textbf{Narrow transformations} (map, filter, flatMap) have each
        output partition depending on at most one input partition and
        require no shuffle. \textbf{Wide transformations} (reduceByKey,
        groupByKey, join) require all-to-all data exchange (shuffle)
        and create stage boundaries.

  \item \textbf{\texttt{reduceByKey}} applies partial aggregation within
        each partition before the shuffle, reducing network traffic by
        approximately the partition count. \textbf{\texttt{groupByKey}}
        transfers all raw values across the network and should be avoided
        whenever an aggregation function is associative and commutative.

  \item The \textbf{Catalyst optimizer} processes DataFrame/SQL plans in
        four phases: analysis (name resolution), logical optimisation
        (predicate pushdown, column pruning, constant folding), physical
        planning (join strategy selection), and code generation (Tungsten
        whole-stage code generation).

  \item \textbf{RDD persistence levels} range from \texttt{MEMORY\_ONLY}
        (fastest, no fault tolerance for evicted partitions) to
        \texttt{DISK\_ONLY} (slowest, full fault tolerance). Kryo
        serialisation reduces the memory footprint and serialisation
        time of persisted RDDs by 3--10$\times$ over Java serialisation.

  \item The \textbf{unified memory model} divides executor heap into
        Reserved (fixed 300\MB), Spark (0.6 fraction), and User
        (0.4 fraction) regions. Within the Spark region, execution and
        storage memory share space dynamically, with execution memory
        able to evict cached RDD partitions.

  \item \textbf{Data skew} occurs when a small number of keys dominate
        the dataset. The \textbf{salting} remedy appends a random integer
        to skewed keys, distributing their records across multiple
        partitions, then removes the salt in a second aggregation pass.

  \item The MLlib \textbf{Pipeline API} expresses preprocessing and
        training as a sequence of Transformers and Estimators. A
        \texttt{Pipeline.fit()} call trains all Estimators and produces
        a \texttt{PipelineModel} that can be saved to HDFS and applied
        to new data.
\end{itemize}

%%====================================================================
\section*{Exercises}
\label{sec:ch6-exercises}
%%====================================================================

\exercisesection{Short Questions}

\sq Explain the MapReduce iteration I/O problem using the formula
$\text{Total I/O} = I \times D \times r$. For a $k$-means clustering
job that runs 50 iterations on a 200\GB\ dataset with $r=3$, calculate
the total disk I/O. How does Spark reduce this?

\sq Define a Resilient Distributed Dataset. List its five defining
properties and explain the purpose of each.

\sq What is lineage-based fault tolerance? Compare it with replication-based
fault tolerance on two dimensions: memory overhead and recovery latency.

\sq Explain the concept of lazy evaluation in Spark. What is the
performance benefit of deferring computation until an action is called?

\sq Distinguish narrow transformations from wide transformations.
Give two examples of each, and explain why wide transformations create
stage boundaries in the DAG.

\sq Why is \texttt{reduceByKey} almost always preferable to
\texttt{groupByKey} for aggregation? Give a concrete numerical example
showing the difference in network traffic for a dataset with 1 billion
records and 200 partitions, where 30\% of records share the same key.

\sq What is a YARN ApplicationMaster for a Spark job? Which component
of Spark architecture runs inside the AM container, and what does it do?

\sq Explain the four phases of the Catalyst optimizer. Which phase is
responsible for choosing between a sort-merge join and a broadcast hash
join, and what criterion determines the choice?

\sq What is the unified memory model in Spark? For an executor with
32\GB\ JVM heap, \texttt{spark.memory.fraction = 0.6}, and
\texttt{spark.memory.storageFraction = 0.5}, calculate the available
storage memory, execution memory, and user memory.

\sq What is Tungsten whole-stage code generation? How does it differ
from the interpreted, row-at-a-time Volcano model used by traditional
query engines?

\sq Explain the Spark MLlib abstraction of Transformer vs.\ Estimator.
Give one concrete example of each and explain what \texttt{fit(df)}
produces for an Estimator.

\sq What is data skew in a Spark wide transformation? Describe the
observable symptom in the Spark UI that indicates a skew problem.

\sq What does \texttt{MEMORY\_ONLY\_SER} mean as an RDD storage level?
Under what circumstances should you choose it over \texttt{MEMORY\_ONLY}?

\sq Explain what a checkpoint does in Spark. When is checkpointing
necessary, and what are its costs?

\sq A Spark job with 8 stages takes 4 hours. The Spark UI shows that
Stage 3 ran for 3.5 hours, with 199 tasks completing in 2 minutes and
one task running for 3.5 hours. What is the most likely cause, and
how would you diagnose it?

\bigskip
\exercisesection{Analytical Problems}

\ap \textbf{Spark vs.\ MapReduce: Iterative Algorithm Cost Analysis.}
A genomics research lab runs a principal component analysis (PCA) on a
SNP (single nucleotide polymorphism) dataset. The dataset is 500\GB,
stored on HDFS with $r = 3$. PCA requires 80 power-iteration steps.

\begin{enumerate}[label=(\alph*)]
  \item Calculate the total disk I/O for a MapReduce implementation
        (one job per iteration, each job reads from and writes to HDFS).
  \item The lab's HDFS cluster sustains 2\GB/s aggregate disk read
        bandwidth. Estimate the wall-clock time for the MapReduce
        implementation.
  \item With Spark and \texttt{MEMORY\_ONLY} persistence, only the
        first iteration reads from HDFS. Subsequent iterations read
        from memory at 50\GB/s aggregate throughput. Estimate the total
        wall-clock time for the Spark implementation.
  \item Calculate the speedup ratio $T_{\text{MR}} / T_{\text{Spark}}$.
  \item The dataset grows to 2\TB. The cluster has 1\TB\ of aggregate
        executor RAM. Explain what happens to the Spark job when the RDD
        no longer fits in memory under \texttt{MEMORY\_ONLY} and under
        \texttt{MEMORY\_AND\_DISK}.
\end{enumerate}

\ap \textbf{Stage and Shuffle Analysis.}
A Spark job reads a log file from HDFS, parses each line into a
(session\_id, page\_url, duration\_seconds) tuple, filters records with
duration $> 5$ seconds, computes the total duration per session, keeps
sessions with total duration $> 60$ seconds, joins with a 200\MB\
session-metadata table (session\_id, user\_id, country), groups by
country, and writes to HDFS.

\begin{enumerate}[label=(\alph*)]
  \item Draw the RDD lineage DAG for this job. Label each RDD with the
        operation that produced it.
  \item Identify all stage boundaries and explain why each boundary
        is placed where it is.
  \item Identify any narrow transformations that can be pipelined into
        a single task pass, and explain the I/O savings this produces.
  \item The session-metadata table is 200\MB. Explain how the Catalyst
        optimizer would handle the join with this table and why this
        avoids a shuffle.
  \item The log file has 20 billion records and \texttt{spark.sql.
        shuffle.partitions = 200}. The join produces 150 billion
        output records. Explain whether 200 shuffle partitions is
        appropriate and how you would determine the correct value.
\end{enumerate}

\ap \textbf{Unified Memory Sizing.}
An e-commerce company runs a Spark job on a YARN cluster. Each executor
is allocated 24\GB\ JVM heap and 6 vCPUs. The job caches a 180\GB\ RDD
(serialised with Kryo) that must be accessed by 20 iterative training
passes. The remaining 30\GB\ of the dataset does not need caching.

\begin{enumerate}[label=(\alph*)]
  \item With \texttt{spark.memory.fraction = 0.6} and
        \texttt{spark.memory.storageFraction = 0.5}, calculate the
        storage memory and execution memory per executor.
  \item If the executor count is 40, calculate the total aggregate
        storage memory across the cluster. Does the 180\GB\ RDD fit?
  \item If Kryo serialisation achieves a 3$\times$ compression ratio
        over the raw RDD size, recalculate whether the RDD fits in
        aggregate storage memory.
  \item If the RDD still does not fit, describe two configuration
        changes (with specific parameter names and values) that could
        increase the available storage memory.
  \item Explain the performance consequence if the RDD does not fit
        in \texttt{MEMORY\_ONLY} storage but the job uses
        \texttt{MEMORY\_AND\_DISK} instead.
\end{enumerate}

\ap \textbf{Data Skew Diagnosis and Salting.}
A retail company runs a Spark job that joins a 5\TB\ transactions table
(partitioned by \texttt{store\_id}) with a 2\TB\ products table. The join
key is \texttt{product\_id}. Analysis of the data reveals:

\begin{center}
\begin{tabular}{lr}
\toprule
\textbf{product\_id} & \textbf{Transaction share} \\
\midrule
\texttt{IPHONE-15-PRO} & 28.3\% \\
\texttt{SAMSUNG-S24}   & 11.7\% \\
\texttt{AIRPODS-PRO}   & 6.2\% \\
All other 4.2M products & 53.8\% \\
\bottomrule
\end{tabular}
\end{center}

\begin{enumerate}[label=(\alph*)]
  \item With 200 shuffle partitions, estimate the number of records
        assigned to the task handling \texttt{IPHONE-15-PRO} vs.\ the
        average task (the total dataset has 50 billion records).
  \item The average task runs in 4 minutes. Estimate the stage
        latency with and without skew mitigation.
  \item Design a salting strategy for the top-3 skewed keys with a
        salt factor of 30. Show the key before and after salting for
        a sample record.
  \item Explain the two-phase aggregation required after salting and
        why it is necessary.
  \item An alternative to salting for join skew is \emph{skew join
        hint} (\texttt{SKEW JOIN} in Spark SQL). Explain conceptually
        how the optimizer would handle a skew hint differently from a
        standard sort-merge join.
\end{enumerate}

\bigskip
\exercisesection{Design Problems}

\dpr \textbf{Spark Architecture for a Real-Time Fraud Detection Platform.}
A global payment network processes 30{,}000 transactions per second.
Each transaction has 45 features (amount, merchant category, cardholder
history, location, time-since-last-transaction, etc.). A fraud detection
model must score each transaction within 200 milliseconds of arrival.
The ML team retrains the model weekly on 90 days of historical
transactions ($\approx$240\TB).

Design the Spark architecture for both the training pipeline and the
serving infrastructure. Your design must specify:

\begin{enumerate}[label=(\alph*)]
  \item The cluster topology for weekly model retraining: number of
        executor nodes, memory per executor, and YARN queue configuration.
        Justify the memory sizing based on the 240\TB\ dataset and the
        10-iteration gradient boosting training algorithm.
  \item The MLlib Pipeline stages for feature engineering: which
        Estimators and Transformers are needed, and in what order.
  \item How the trained model is exported from the Spark cluster to
        the serving layer, given the 200-millisecond latency constraint
        (Spark jobs take seconds to initialise; model serving cannot
        invoke a Spark job per transaction).
  \item The persistence strategy for the feature engineering pipeline:
        which intermediate RDDs should be cached, at what storage level,
        and why.
  \item How the design changes if the retraining frequency is increased
        from weekly to hourly, incorporating the most recent 6 hours of
        transactions into the model on each run.
\end{enumerate}

\dpr \textbf{Spark SQL Migration from Hive-on-MapReduce.}
A telecommunications company's analytics team runs 400 Hive-on-MapReduce
queries per day against a 500\TB\ ORC warehouse. Average query latency
is 45 minutes; the SLA for analyst-facing queries is 15 minutes; and
50 queries (12.5\%) exceed 2 hours. The data engineering team proposes
migrating to Spark SQL.

Design the migration. Your design must address:

\begin{enumerate}[label=(\alph*)]
  \item The Catalyst optimisations (predicate pushdown, column pruning,
        broadcast joins) that are most likely to reduce latency for
        queries over an ORC table partitioned by date and region.
        Estimate the I/O reduction for a query selecting 8 of 120
        columns, filtering on a date range covering 5\% of the total
        data, with an additional predicate that eliminates 70\% of the
        remaining row groups.
  \item The YARN queue configuration changes required to co-run Hive-on-MapReduce
        (existing nightly ETL) and Spark SQL (new analyst workload) on the
        same cluster without either workload starving the other.
  \item A compatibility verification strategy: how to confirm that the
        Spark SQL query results match the Hive-on-MapReduce results for
        the 400 queries before switching production traffic.
  \item The expected performance change for the 50 slow queries (currently
        $>$2 hours). Which of the following root causes would be fixed
        by Catalyst, and which would still require manual query rewriting:
        (i) full-table scans that could use partition pruning;
        (ii) \texttt{GROUP BY} on a column with extreme skew;
        (iii) a three-table join where the smallest table is 400\MB.
  \item A rollback plan if Spark SQL produces incorrect results for
        a category of queries (e.g., NULL handling differences from Hive's
        default behaviour).
\end{enumerate}

\bigskip
\exercisesection{Programming Exercises}

\pe \textbf{Word Frequency with RDD API (PySpark).}
Implement a word frequency counter using the Spark RDD API. The program
must: (1) read a text file from HDFS; (2) tokenise each line into words
(lowercase, strip punctuation); (3) count occurrences of each word; (4)
return the top 20 most frequent words and their counts. Measure the
performance difference between using \texttt{reduceByKey} and
\texttt{groupByKey}.

\begin{lstlisting}[style=pyspark, caption={PE6.1: Word frequency with RDD API}]
import re
from pyspark.sql import SparkSession

spark = SparkSession.builder.appName("WordFreq").getOrCreate()
sc    = spark.sparkContext

def tokenize(line):
    """Lower-case and split on non-alpha characters."""
    return [w for w in re.split(r"[^a-z]+", line.lower()) if w]

# Read text from HDFS (replace with actual path)
text_rdd = sc.textFile("hdfs:///data/wikipedia_sample.txt")

# --- Implementation using reduceByKey (correct approach) ---
freq_rdd = (text_rdd
    .flatMap(tokenize)
    .map(lambda w: (w, 1))
    .reduceByKey(lambda a, b: a + b)
    .sortBy(lambda kv: kv[1], ascending=False)
)

top20 = freq_rdd.take(20)
print("Top 20 words (reduceByKey):")
for word, count in top20:
    print(f"  {word:<20} {count:>8,}")

# --- Implementation using groupByKey (anti-pattern for comparison) ---
# TODO: implement using groupByKey and compare memory/time
# groupby_rdd = (text_rdd
#     .flatMap(tokenize)
#     .groupBy(lambda w: w)
#     .mapValues(lambda vals: sum(1 for _ in vals))
#     ...
# )

# Cache and measure second-pass performance
freq_rdd.persist()
print(f"\nTotal unique words: {freq_rdd.count():,}")
\end{lstlisting}

\pe \textbf{Iterative $k$-Means with RDD Caching (PySpark).}
Implement $k$-means clustering using the raw RDD API (not MLlib) to
observe the effect of caching on iterative performance. The implementation
must: (1) cache the input data RDD after the first iteration; (2) measure
wall-clock time per iteration; (3) report the speedup of cached vs.\
uncached iterations.

\begin{lstlisting}[style=pyspark, caption={PE6.2: Iterative k-means with RDD caching}]
import time
import random
import math
from pyspark.sql import SparkSession
from pyspark import StorageLevel

spark = SparkSession.builder.appName("KMeans").getOrCreate()
sc    = spark.sparkContext

def parse_point(line):
    return [float(x) for x in line.strip().split(",")]

def distance(p, c):
    return math.sqrt(sum((a - b)**2 for a, b in zip(p, c)))

def closest_centroid(point, centroids):
    dists = [distance(point, c) for c in centroids]
    return dists.index(min(dists))

def add_vectors(v1, v2):
    return [a + b for a, b in zip(v1, v2)]

def mean_vector(total, count):
    return [x / count for x in total]

K         = 10
MAX_ITERS = 20
SEED      = 42

# Load data
data = sc.textFile("hdfs:///data/clustering_input.csv").map(parse_point)

# Cache the input data  --  this is the key optimisation
data.persist(StorageLevel.MEMORY_ONLY)
data.count()    # trigger caching (materialize into memory)

# Initialise centroids randomly from the first partition
centroids = data.takeSample(False, K, seed=SEED)

iter_times = []
for iteration in range(MAX_ITERS):
    t_start = time.time()

    # Assign each point to its nearest centroid
    assignments = data.map(
        lambda p: (closest_centroid(p, centroids),
                   (p, 1))
    )

    # Sum vectors and counts per cluster
    cluster_stats = (assignments
        .reduceByKey(lambda a, b: (add_vectors(a[0], b[0]),
                                   a[1] + b[1]))
    )

    # Compute new centroids
    new_centroids_list = (cluster_stats
        .mapValues(lambda vs: mean_vector(vs[0], vs[1]))
        .collect()
    )
    new_centroids = {k: v for k, v in new_centroids_list}
    centroids = [new_centroids.get(i, centroids[i]) for i in range(K)]

    elapsed = time.time() - t_start
    iter_times.append(elapsed)
    print(f"Iteration {iteration+1:>2}: {elapsed:.2f}s")

print(f"\nFirst iteration : {iter_times[0]:.2f}s  (HDFS read + compute)")
print(f"Mean iter 2-20  : {sum(iter_times[1:])/len(iter_times[1:]):.2f}s  (cached)")
print(f"Speedup from cache: {iter_times[0]/sum(iter_times[1:])*len(iter_times[1:]):.1f}x")
\end{lstlisting}

\pe \textbf{Salting Benchmark for Skew Mitigation (PySpark).}
Simulate a skewed dataset and measure the performance difference between
a naive \texttt{reduceByKey} (with skew) and a salted two-phase
aggregation. Report per-task statistics to demonstrate that salting
eliminates the long tail.

\begin{lstlisting}[style=pyspark, caption={PE6.3: Salting benchmark for skew mitigation}]
import random
import time
from pyspark.sql import SparkSession

spark = SparkSession.builder.appName("SkewBenchmark") \
    .config("spark.sql.shuffle.partitions", "50") \
    .getOrCreate()
sc = spark.sparkContext

TOTAL_RECORDS  = 5_000_000
SKEW_KEY       = "IPHONE-15-PRO"
SKEW_FRACTION  = 0.35           # 35% of all records
SALT_FACTOR    = 20             # number of salt buckets

def generate_record(_):
    """Generate (product_id, sale_amount) with skew."""
    if random.random() < SKEW_FRACTION:
        return (SKEW_KEY, random.randint(500, 1500))
    products = [f"PROD-{i:06d}" for i in range(1, 1001)]
    return (random.choice(products), random.randint(10, 500))

data = sc.parallelize(range(TOTAL_RECORDS), 50).map(generate_record)
data.persist()
data.count()    # materialise

# --- Approach 1: naive reduceByKey (skewed) ---
t0 = time.time()
naive_result = data.reduceByKey(lambda a, b: a + b).collect()
t_naive = time.time() - t0
print(f"Naive reduceByKey: {t_naive:.2f}s")

# --- Approach 2: salted two-phase aggregation ---
SKEWED_KEYS = {SKEW_KEY}

def salt_key(record):
    key, val = record
    if key in SKEWED_KEYS:
        return ((key, random.randint(0, SALT_FACTOR - 1)), val)
    return ((key, 0), val)

def desalt_key(record):
    (key, _salt), val = record
    return (key, val)

t1 = time.time()
salted_result = (data
    .map(salt_key)
    .reduceByKey(lambda a, b: a + b)
    .map(desalt_key)
    .reduceByKey(lambda a, b: a + b)
    .collect()
)
t_salted = time.time() - t1
print(f"Salted 2-phase  : {t_salted:.2f}s")

print(f"Speedup         : {t_naive/t_salted:.2f}x")

# Verify correctness: totals should match
naive_dict  = dict(naive_result)
salted_dict = dict(salted_result)
assert abs(naive_dict.get(SKEW_KEY, 0) -
           salted_dict.get(SKEW_KEY, 0)) < 1e-6, "Results differ!"
print("Correctness check: PASSED")
\end{lstlisting}
